\documentclass[11pt, a4paper, english]{scrartcl}

% --- Language & Encoding ---
\usepackage[utf8]{inputenc}
\usepackage[T1]{fontenc}
\usepackage{babel}

% --- Fonts & Typography ---
\usepackage{charter}
\usepackage[scaled=0.92]{helvet}
\usepackage{inconsolata}
\usepackage{microtype}

% --- Page Geometry ---
\usepackage[
  top=2.2cm,
  bottom=2.5cm,
  left=2.3cm,
  right=2.3cm,
  headheight=14pt
]{geometry}

% --- Colors ---
\usepackage[table]{xcolor}
\definecolor{NavyBlue}{RGB}{20,40,75}
\definecolor{TealAccent}{RGB}{0,120,130}
\definecolor{WarmGold}{RGB}{185,125,40}
\definecolor{DarkSlate}{RGB}{45,55,72}
\definecolor{LightBg}{RGB}{246,248,251}
\definecolor{BorderGrey}{RGB}{220,225,232}
\definecolor{AlertRed}{RGB}{160,40,40}

% --- Headers & Footers ---
\usepackage{scrlayer-scrpage}
\pagestyle{scrheadings}
\clearscrheadings
\ihead{\footnotesize\sffamily\color{DarkSlate}CS 8803: Advanced Distributed Systems \textbar{} Seminar Response}
\ohead{\footnotesize\sffamily\color{DarkSlate}Week 7: Fault-Tolerant Consensus}
\cfoot{\pagemark}
\addtokomafont{pagefoot}{\small\sffamily\color{DarkSlate}}
\KOMAoptions{headsepline=0.5pt}
\addtokomafont{headsepline}{\color{BorderGrey}}

% --- Math & Tables ---
\usepackage{amsmath,amssymb}
\usepackage{booktabs}
\usepackage{tabularx}
\definecolor{TableHeadBg}{RGB}{235,240,248}
\usepackage{enumitem}
\usepackage{graphicx}
\usepackage{fontawesome5}

% --- TikZ & Graphics ---
\usepackage{tikz}
\usetikzlibrary{calc,positioning,shapes.geometric}

% --- tcolorbox ---
\usepackage[many]{tcolorbox}
\tcbuselibrary{skins,breakable}

% --- Hyperlinks ---
\usepackage[
  colorlinks=true,
  linkcolor=TealAccent,
  citecolor=TealAccent,
  urlcolor=TealAccent
]{hyperref}

% --- Section Headings Styling (KOMA-Script Native) ---
\setkomafont{section}{\normalfont\Large\sffamily\bfseries\color{NavyBlue}}
\setkomafont{subsection}{\normalfont\large\sffamily\bfseries\color{TealAccent}}

\makeatletter
\renewcommand*{\sectionlinesformat}[4]{%
  \@hangfrom{\hskip #2#3}{#4}%
  \Ifstr{#1}{section}{\par\nobreak\vspace{2pt}{\color{TealAccent}\hrule height 0.8pt}}{}%
}
\makeatother

% --- Custom tcolorbox Styles ---
\tcbset{
  paperbox/.style={
    enhanced,
    colback=LightBg,
    colframe=NavyBlue,
    arc=3pt,
    boxrule=0.8pt,
    left=10pt,
    right=10pt,
    top=8pt,
    bottom=8pt,
    fonttitle=\bfseries\sffamily\color{white},
    coltitle=white,
    colbacktitle=NavyBlue,
    attach boxed title to top left={xshift=10pt, yshift=-10pt},
    boxed title style={arc=2pt, boxrule=0pt, top=4pt, bottom=4pt, left=8pt, right=8pt}
  },
  promptbox/.style={
    enhanced,
    colback=white,
    colframe=TealAccent,
    arc=3pt,
    boxrule=1pt,
    left=12pt,
    right=12pt,
    top=8pt,
    bottom=8pt,
    drop shadow=black!10
  },
  takeawaybox/.style={
    enhanced,
    colback=LightBg,
    colframe=WarmGold,
    arc=2pt,
    boxrule=0.6pt,
    borderline west={3pt}{0pt}{WarmGold},
    left=10pt,
    right=10pt,
    top=8pt,
    bottom=8pt
  }
}

\begin{document}

% ==================== HEADER BLOCK ====================
\begin{tcolorbox}[
  enhanced,
  colback=NavyBlue,
  colframe=NavyBlue,
  arc=4pt,
  top=14pt,
  bottom=14pt,
  left=16pt,
  right=16pt
]
\begin{minipage}[t]{0.62\textwidth}
  {\color{WarmGold}\sffamily\bfseries\footnotesize GRADUATE SEMINAR READING RESPONSE \textbar{} WEEK 7}\\[4pt]
  {\color{white}\sffamily\huge\bfseries High-Latency Fault Tolerance \& Consensus Protocols}\\[6pt]
  {\color{LightBg}\sffamily\small CS 8803: Advanced Distributed Infrastructure \& Network Resilience}
\end{minipage}%
\hfill
\begin{minipage}[t]{0.34\textwidth}
  \color{white}\sffamily\small
  \begin{tabular}{rl}
    \textbf{Student:} & Alex Sterling \\
    \textbf{NID:} & N-19842091 \\
    \textbf{Instructor:} & Prof. Evelyn Vance \\
    \textbf{Date:} & October 24, 2026 \\
    \textbf{Track:} & Distributed Systems
  \end{tabular}
\end{minipage}
\end{tcolorbox}

\vspace{0.5em}

% ==================== SECTION 1: ASSIGNED READINGS ====================
\section{Assigned Readings \& Metadata}

\begin{tcolorbox}[paperbox, title={\faBook\ \ Reading Portfolio}]
\small
\begin{tabularx}{\linewidth}{@{}l X l l@{}}
  \textbf{Ref} & \textbf{Title \& Authors} & \textbf{Venue / Year} & \textbf{Primary Focus} \\
  \toprule
  \textbf{\lbrack P1\rbrack} & \textbf{Speculative State Replication in Partial Asynchrony}\newline \emph{Marcus Thorne, Elena Rostova, and Julian Sterling} & ACM TOCS (2024) & Consensus Latency \\
  \midrule
  \textbf{\lbrack P2\rbrack} & \textbf{Dynamic Quorum Reconfiguration under Byzantine Churn}\newline \emph{Sarah Jenkins and David K. Chen} & OSDI '25 & Adaptive Membership \\
  \midrule
  \textbf{\lbrack P3\rbrack} & \textbf{Practical Bounded Memory Footprints for Geo-Raft}\newline \emph{Nexa Research Group (Tech Report \#402)} & Tech Report (2026) & Storage Overhead \\
  \bottomrule
\end{tabularx}
\end{tcolorbox}

\begin{tcolorbox}[takeawaybox]
\sffamily\small
\textbf{\faLightbulb\ Executive Summary of Week 7 Scope:}
This week's readings examine the tension between network latency and safety guarantees in fault-tolerant replicated state machines. \textbf{\lbrack P1\rbrack} introduces optimistic fast-path execution reducing round-trips from $3 \rightarrow 1$ under zero-contention scenarios. \textbf{\lbrack P2\rbrack} relaxes static node set assumptions by formalizing epoch-based dynamic reconfiguration under active adversarial churn. \textbf{\lbrack P3\rbrack} addresses log compaction memory explosion in multi-region deployments managed by Synapse Infrastructure Engine.
\end{tcolorbox}

% ==================== SECTION 2: SYNTHESIS & COMPARATIVE ANALYSIS ====================
\section{Core Synthesis \& Comparative Analysis}

\subsection{Architectural Paradigms: Optimistic vs. Pessimistic Consensus}

The central thesis connecting \textbf{\lbrack P1\rbrack} and \textbf{\lbrack P2\rbrack} revolves around the tradeoff between fast-path speculation and worst-case protocol fallback. Traditional Multi-Paxos and Raft enforce structural ordering prior to state application. Thorne et al. (\textbf{\lbrack P1\rbrack}) break this barrier by allowing replicas to speculatively execute incoming request batches provided a supermajority ($Q_{\text{fast}} = \lceil \frac{3f+1}{2} \rceil$) of nodes agree on the sequence vector.

Mathematically, while classical PBFT requires $3f+1$ total nodes to tolerate $f$ Byzantine faults with 2-phase commit (2 message round-trips), \textbf{\lbrack P1\rbrack} achieves single round-trip consensus whenever the network latency ratio satisfies:
\begin{equation}
  \Delta_{\text{actual}} \le \frac{1}{2} \delta_{\text{estimated}} + \gamma_{\text{jitter}}
\end{equation}
where $\gamma_{\text{jitter}}$ represents the variance in inter-replica packet transit time across WAN links.

\begin{table}[h]
\centering
\caption{\color{NavyBlue}\sffamily Comparative Matrix of Protocol Guarantees}
\vspace{4pt}
\small
\begin{tabularx}{\linewidth}{@{} l X X X @{}}
  \toprule
  \rowcolor{TableHeadBg}
  \textbf{Dimension} & \textbf{Thorne et al. \lbrack P1\rbrack} & \textbf{Jenkins \& Chen \lbrack P2\rbrack} & \textbf{Nexa Group \lbrack P3\rbrack} \\
  \midrule
  \textbf{Network Assumption} & Partial Asynchrony ($\text{GST}$) & Asynchronous WAN & Synchronous LAN / Edge \\
  \textbf{Fault Model} & Crash-Recovery + Byzantine & Byzantine Churn ($f < n/3$) & Fail-Stop Only \\
  \textbf{Commit Latency} & 1 RTT (Optimistic) / 3 RTT & 2 RTT + Epoch Sync & 2 RTT (Direct Log Stream) \\
  \textbf{Reconfiguration} & Static Configuration & Epoch-based Dynamic Quorums & Joint Consensus (Raft) \\
  \textbf{Memory Bound} & $O(N \cdot \text{Window})$ & $O(N \cdot \text{Epochs})$ & $O(\text{Active Snapshots})$ \\
  \bottomrule
\end{tabularx}
\end{table}

\subsection{Trade-offs in Dynamic Quorum Maintenance}

Jenkins and Chen (\textbf{\lbrack P2\rbrack}) pick up where \textbf{\lbrack P1\rbrack} leaves off by addressing node departure and joining in WAN environments. While \textbf{\lbrack P1\rbrack} assumes a static universe of $N$ processes, real-world cloud infrastructures (such as Vertex Distributed Storage or Apex Engine) undergo constant maintenance, partition recovery, and spot instance terminations.

The authors propose \emph{Epoch-Shift}, a mechanism that decouples state machine execution from membership agreement. Instead of stalling the consensus pipeline during reconfiguration, \textbf{\lbrack P2\rbrack} runs a lightweight parallel consensus instance to elect the epoch configuration $E_{k+1}$ while $E_k$ processes client payload writes. However, this approach exposes an interesting edge case: if $E_k$ and $E_{k+1}$ experience concurrent partitions, client requests can be committed in $E_k$ but omitted from the initial snapshot transferred to $E_{k+1}$.

% ==================== SECTION 3: CRITICAL EVALUATION & WEAKNESSES ====================
\section{Critical Evaluation \& Methodological Weaknesses}

\begin{enumerate}[label={\sffamily\bfseries W\arabic*.}, leftmargin=2em]
  \item \textbf{Unrealistic Asymmetry Assumptions in Thorne et al. \lbrack P1\rbrack:}\\{}
  The paper evaluates speculative state replication over a controlled 5-region AWS deployment. However, their synthetic workload benchmark assumes symmetric link latencies between all leader-follower pairs. In production networks managed by Nimbus Cloud, cross-Atlantic latency exhibits significant jitter ($\sigma^2 > 45\text{ms}$). Under such variance, the fast-path condition ($Q_{\text{fast}}$) falls back to the slow-path in $38\%$ of trials, incurring a double-penalty of execution rollback and state repair.

  \item \textbf{Unbounded State Growth during Reconfiguration Catch-up \lbrack P2\rbrack:}\\{}
  Jenkins and Chen rely on sliding-window state checkpoints to onboard new quorum members in \textbf{\lbrack P2\rbrack}. When churn rates exceed $12\%$ per epoch, the delta log required to synchronize a joining node exceeds the snapshot buffer size. The paper fails to quantify the threshold at which memory exhaustion triggers cascading node evictions.

  \item \textbf{Over-Reliance on Hardware Security Modules (HSM) in Nexa Tech Report \lbrack P3\rbrack:}\\{}
  \textbf{\lbrack P3\rbrack} achieves impressive memory compaction numbers ($74\%$ reduction in RAM consumption) by offloading signature verification and log sequence hashing to dedicated TPM/HSM units. This assumption limits applicability in modern serverless environments (e.g., Meridian Edge Functions) where hardware enclave access is restricted or virtualization overhead negates throughput gains.
\end{enumerate}

% ==================== SECTION 4: SEMINAR DISCUSSION PROMPTS ====================
\section{Discussion Questions \& Seminar Prompts}

\begin{tcolorbox}[promptbox]
\sffamily
\textbf{\color{NavyBlue}\faQuestionCircle\ Prompt 1: Speculative Rollbacks vs. Storage Durability}
\smallskip

In \textbf{\lbrack P1\rbrack}, fast-path speculative execution modifies memory state before disk synchronization completes. If a power failure occurs simultaneously across $f+1$ nodes during an uncommitted speculative window, can state corruption be strictly isolated without invalidating non-volatile RAM logs? How would you design a recovery protocol that avoids replaying phantom operations?

\begin{itemize}[label={\color{TealAccent}\small\faAngleRight}, leftmargin=1.2em, itemsep=2pt]
  \item \emph{Key Angles:} NVRAM persistence guarantees, linearizability vs. sequential consistency, hardware costs.
  \item \emph{Target Audience:} Distributed storage engineers and transaction processing specialists.
\end{itemize}
\end{tcolorbox}

\vspace{0.4em}

\begin{tcolorbox}[promptbox]
\sffamily
\textbf{\color{NavyBlue}\faQuestionCircle\ Prompt 2: Governance \& Security under High Churn Rates}
\smallskip

Jenkins and Chen (\textbf{\lbrack P2\rbrack}) argue that epoch transitions can occur as frequently as every 5 seconds under severe network volatility. Does this high re-configuration frequency open a novel attack vector where a malicious actor can orchestrate a Sybil attack by selectively delaying epoch ACK messages, thereby freezing the consensus membership indefinitely?

\begin{itemize}[label={\color{TealAccent}\small\faAngleRight}, leftmargin=1.2em, itemsep=2pt]
  \item \emph{Key Angles:} Adversarial network manipulation, protocol liveness vs. safety under active Denial of Service.
  \item \emph{Target Audience:} Fault-tolerant protocol designers and security researchers.
\end{itemize}
\end{tcolorbox}

\vspace{0.4em}

\begin{tcolorbox}[promptbox]
\sffamily
\textbf{\color{NavyBlue}\faQuestionCircle\ Prompt 3: Real-World Viability of Hybrid Consensus Engines}
\smallskip

Comparing \textbf{\lbrack P1\rbrack} and \textbf{\lbrack P3\rbrack}, is it practical to construct a hybrid state machine engine that dynamically switches between speculative fast-paths during peak hours and compact storage-constrained Raft modes during batch maintenance windows? What metadata overhead would such runtime mode switching introduce?

\begin{itemize}[label={\color{TealAccent}\small\faAngleRight}, leftmargin=1.2em, itemsep=2pt]
  \item \emph{Key Angles:} Operational complexity, dynamic reconfiguration bugs, production maintainability.
\end{itemize}
\end{tcolorbox}

% ==================== SECTION 5: FUTURE RESEARCH DIRECTIONS ====================
\section{Synthesis of Future Research Directions}

Based on this week's papers, three promising research avenues emerge for potential term projects or thesis work:

\begin{figure}[h]
\centering
\begin{tikzpicture}[node distance=0.8cm, auto, >=stealth]
  \node [draw=NavyBlue, fill=LightBg, rounded corners=3pt, inner sep=6pt, text width=3.6cm, align=center, font=\sffamily\small] (spec)
    {\textbf{Speculative WAN Consensus}\\(\textbf{\lbrack P1\rbrack} Extension)};
  \node [draw=TealAccent, fill=LightBg, rounded corners=3pt, inner sep=6pt, text width=3.6cm, align=center, font=\sffamily\small, right=0.6cm of spec] (dynamic)
    {\textbf{Zero-Trust Dynamic Quorums}\\(\textbf{\lbrack P2\rbrack} Extension)};
  \node [draw=WarmGold, fill=LightBg, rounded corners=3pt, inner sep=6pt, text width=3.6cm, align=center, font=\sffamily\small, right=0.6cm of dynamic] (compaction)
    {\textbf{Hardware-Assisted Compaction}\\(\textbf{\lbrack P3\rbrack} Extension)};

  \draw[->, thick, NavyBlue] (spec) -- node[above, font=\sffamily\tiny\color{DarkSlate}]{Epoch Handshake} (dynamic);
  \draw[->, thick, TealAccent] (dynamic) -- node[above, font=\sffamily\tiny\color{DarkSlate}]{Log Pruning} (compaction);
\end{tikzpicture}
\caption{\color{NavyBlue}\sffamily Interconnection of Future Research Tracks}
\end{figure}

\begin{description}[font=\sffamily\bfseries\color{TealAccent}, leftmargin=1.5em]
  \item[Track A: Machine-Learned Latency Predictors for Speculative Paths]
  Rather than using static threshold estimates $\delta_{\text{estimated}}$ as in \textbf{\lbrack P1\rbrack}, replace the fast-path trigger with a lightweight online neural predictor that forecasts link congestion based on real-time telemetry from Vertex edge routers.

  \item[Track B: Formally Verified Reconfiguration in TLA+]
  Neither \textbf{\lbrack P1\rbrack} nor \textbf{\lbrack P2\rbrack} provides mechanically verified safety proofs for concurrent epoch updates during split-brain recovery. A valuable contribution would be modeling \textbf{\lbrack P2\rbrack}'s Epoch-Shift protocol in TLA+ / TLC to verify invariant retention under network partitions.
\end{description}

% ==================== REFERENCES ====================
\section{References \& Extended Bibliography}

\small
\begin{enumerate}[label={\sffamily\bfseries\color{NavyBlue}[\arabic*]}, leftmargin=2.2em, itemsep=3pt]
  \item Thorne, M., Rostova, E., \& Sterling, J. (2024). Speculative State Replication in Partial Asynchrony. \emph{ACM Transactions on Computer Systems (TOCS)}, 42(3), 112--138.
  \item Jenkins, S., \& Chen, D. K. (2025). Dynamic Quorum Reconfiguration under Byzantine Churn. In \emph{Proc. 19th USENIX Symposium on Operating Systems Design and Implementation (OSDI '25)}, pp. 405--422.
  \item Nexa Research Group (2026). Practical Bounded Memory Footprints for Geo-Replicated Raft Systems. \emph{Technical Report \#402}, Nexa Institute of Distributed Infrastructure.
  \item Lamport, L. (1998). The Part-Time Parliament. \emph{ACM Transactions on Computer Systems}, 16(2), 133--169.
  \item Ongaro, D., \& Ousterhout, J. (2014). In Search of an Understandable Consensus Algorithm. In \emph{Proc. USENIX Annual Technical Conference (ATC '14)}, pp. 305--319.
\end{enumerate}

\end{document}
